\documentclass[11pt]{article}

%%| INCLUDE PACKAGES HERE |%%
\usepackage[utf8]{inputenc}
\usepackage[paper=a4paper,left=25mm,right=25mm,top=25mm,bottom=25mm]{geometry}
\usepackage[english, ngerman]{babel}
\usepackage{graphicx}
\usepackage{amsmath}
\usepackage{amssymb}
\usepackage[autostyle=true,german=quotes]{csquotes}
\usepackage[backend=biber, style=ieee]{biblatex}
\usepackage{hyperref}
\usepackage{setspace}
\usepackage{microtype}

%%| SET PACKAGE PARAMETERS HERE |%%
\hypersetup{                                                                        
    colorlinks,
    citecolor   = black,
    filecolor   = black,
    linkcolor   = black,
    urlcolor    = black,
    pdftitle    = {Labor ABCDEFGHIJ},                                               % SET FILE PARAMETERS HERE (SHOWN IN PDF PROPERTIES)
    pdfsubject  = {Einführung in ABCDEFGHIJ},                                       % ...
    pdfauthor   = {Max Mustermann},                                                 % ...
    pdfkeywords = {XY-Labor, Laborbericht ...} ,                                    % ...
    pdfcreator  = {pdflatex},           
    pdfproducer = {LaTeX with hyperref}
}
\setstretch{1.25}                                                                   % SET GENERAL LINE SPACING HERE

%%| SET OTHER PARAMETERS HERE |%%
\addbibresource{Document/Bibliography/ref.bib}

%%%%%%%%%%%%%%%%%%%%%%%%%%%%%%%%%%%%%%%%%%%%%%%%%%%%%%%%%%%%%%%%%%%%%%%%%%%%%%%%%%%%%%%%%%%%%%%%%%%%%%%%%%%%%%%%%%%%%%%%%%%%%%%%%%%%%%%
%%| TITLE PAGE |%%
\begin{document}
\begin{titlepage}
    \begin{center}
        \includegraphics[width=0.55\textwidth]{Images/Template/HsKa_new.png}\\[16ex]
        \huge{\textbf{Ethik für Informatik}}\\[8ex]
        \large{\textbf{Take-Home-Aufgabe}}\\[1ex]
        \normalsize{}
        \begin{tabular}{lll}
            Student:            & \quad Erlind Sejdiu                                & 106357 \\[2ex]     % NUMBER IS MAT.NR
            Universität:        & \quad Hochschule Karlsruhe Technik und Wirtschaft   &       \\[2ex]
            Studiengang:        & \quad Informatik- und Wirtschaftsinformatik           &       \\[2ex]
            Studienvertiefung:  & \quad Informatik Master                         &       \\[2ex]
            Semester:           & \quad Sommersemester 2026                             &       \\[2ex]
            Dozent:             & \quad Prof. Dr. Ziad Mahayni                       &       \\[2ex]
            Bearbeitet am:      & \quad 28. Juni 2026                                  &       \\[2ex]
        \end{tabular}
    \end{center}
\end{titlepage}
\newpage
%%%%%%%%%%%%%%%%%%%%%%%%%%%%%%%%%%%%%%%%%%%%%%%%%%%%%%%%%%%%%%%%%%%%%%%%%%%%%%%%%%%%%%%%%%%%%%%%%%%%%%%%%%%%%%%%%%%%%%%%%%%%%%%%%%%%%%%
%%| TABLE OF CONTENTS |%%
\pagenumbering{gobble}    
\tableofcontents
\newpage
\pagenumbering{arabic}

\section{Verantwortungsdilemma im Digitalenzeitalter}
Die digitale Ethik beschäftigt sich mit moralischen Fragen von der Nutzung digitaler Techniken. Dabei hat sich zum Beispiel auch das Verantwortungsdilemma ergeben, aber um dieses Dilema zu verstehen werden zunächst 2 Paradigmen eingeführt.

\paragraph{Techschnische Realisierbarkeit}
Eine festellende Beobachtung ist das immer schnellerer Foranschreiten der Technik. Dabei entsteht ein Paradigma, welches besagt, dass mit neu erforschten Technik sich  der Handlungsspielraum der Menschen erweitert. Als Beispiel hier können die Lebensumstände in der Vergangenheit angeführt werden. Früher ist es typisch gewesen, dass die Menschen nah beinander Lebten, wenn man sich sehen wollte. Heutzutage ist es möglich mit Menschen zu kommunizieren, welche mehrere hundert Kilometer entfernt sind und das aus sehr offentsichlichen Gründen. Wenn ich meine Freunde in der Schule sehen möchte, muss ich nicht mehr zu ihnen fahren, sondern kann einfach mit ihnen über das Internet kommunizieren. Wenn ich diese doch treffen möchte kann ich mit meinem Auto einfach zu ihnen hinfahren. Diese Möglichkeiten waren noch vor 200 jahren stark eingeschränkt. Mit steigenem technologischen Fortschritt erweitert sich die Handlungsmöglichkeiten der Menschen. Je mehr ich in der Lage bin, Aktionen  auszuführen, desto mehr Verantwortung habe ich, diese Dinge auch richtig zu tun \cite{article}.

\paragraph{Technologieadaption}
Eine weiterer Apspekt ist die Adaptionsgeschwindigkeit der Menschen. Dieses Paradigma besagt, dass die Anpassungfähigkeit an eine Technologie langsamer ist als die Technologieentwicklung. In der Vergangenheit war der wechsel von geozentrischen Weltbild zum heliozentrischen Weltbild ein technischer Fortschritt, der anfangs stark verneint worden ist. Die Veränderung, dass die Welt sich um die Sonne dreht wurde am Anfang langsam wahrgenommen. Um genau zu sein hat dieses Paradigmawechsel 150 Jahre gedauert. Im starken Kontrast dazu ist die Digitalisierung, die in den vergangenen Jahren stattgefunden hat und damit neue Technologien erfunden wurden, die in viel kürzerer Zeitspanne stattgefunden hat und wahrgenommen wurden. Das heißt neben der Erfindung neuer Technik und der tatsächlichen Nutzung dieser Technik wird eine gewissen adaptionszeit benötigt \cite{article}. 

\paragraph{Verantwortungsdilemma}

Aus diesen zwei Paradigmen ergibt sicht ein Verantwortungsdilemma. Neue Technologien werden schneller in unseren Alltag integriert, aber weil die Adaptionsgeschwindigkeit der Menschen nicht mit der Geschwindigkeit des technologischen Fortschritts mithalten kann, kann diese neu entstandene Technologie im bestimmten Kontext nicht verstanden werden oder auch missbraucht werden. Dies ist in dem Sinne drastisch, da heutzutage mit Technologien vieles umsetzen können. Ein Beispiel dafür ist die Deepfake-Technologie. Deepfakes sind synthetische Medieninhalte, die mithilfe von künstlicher Intelligenz erstellt werden und es ermöglichen, realistische Videos oder Audiodateien zu erzeugen, in denen Personen Dinge sagen oder tun, die sie tatsächlich nicht gesagt oder getan haben. Diese Technologie hat das Potenzial, die Art und Weise, wie Informationen verbreitet werden, grundlegend zu verändern. Wenn Menschen nicht in der Lage sind, zwischen echten und gefälschten Inhalten zu unterscheiden, kann dies zu einer Infokalypse führen, in der die Verbreitung von Desinformationen und Falschmeldungen außer Kontrolle gerät. Da diese Technologie so stark vorangeschritten ist, sind Regularien und Gesetze, die den Umgang mit dieser Technologie regeln, nicht mehr ausreichend, um die Verbreitung von Deepfakes zu kontrollieren \cite{article}.

\section{Datenethik: AI-Bias}

Die Philosophie lässt sich in mehrere Kategorien unterteilen. Eine dieser Kategorien ist die Datenethik, welche sich in weitere Unterkategorien aufteilt. Diese Unterkategorien sind jeweils AI-Bias, digitale Überwachung, Roboterethik und virtuelle Akteure sowie Maschinenethik \& „Alignment-Problem“. Im Folgenden wird die Datenethik bezüglich der AI-Bias genauer beleuchtet.

\enquote{Dieser Begriff beschreibt die genannte, durch Verzerrungen in den Daten ausgelöste Ungleichbehandlung durch KI-Systeme}\cite[S. ~5]{Bias}. Das heißt, dass bestimmte Personen oder Personengruppen systematisch benachteiligt werden. Dabei ist es erstmal wichtig zu verstehen, dass ein AI-Bias durch unterschiedliche Gründe entstehen kann. Dabei kann der AI-Bias in der Datenerhebung und Präparation, in der Modellentwicklung und Training und im Deployment und im Monitoring stattfinden \cite{Bias}.

Ein Bias bei der Datenerhebung kann durch die Verwendung veralteter Datensätze enstehen, welche Stereotypen wiedergeben. Als konkretes Beispiel kann hier das Stereotyp genommen werden, welches Männer als Arbeiter und Frauen als Hausfrauen darstellt. Wird ein KI-Model mit solchen Daten trainiert, so kann eine Frau bei Bewerbungen systematischer schlechter abschneiden, da das Model der Aufassung ist, dass Männer die Arbeiter sind während Frauen nur Hausfrauen sind. Im Jahr 2017 ist genau so ein Fall in einer Software von Amazon passiert. Diese Software hat Menschen abhängig von Ihren qualifikationen vorselektiert und dabei wurden Frauen in Programmierberufen schlechter bewertet, weil das Model auf einem Datensatz trainiert wurde, welcher großtenteils mit Männdern gefüllt war \cite{dastin_insight_2018, Bias}.  

Auch während der Entwicklung kann ein Bias enstehen. Wird während einer Entwicklung ein Datensatz verwendet werden, welcher Gesichter aus bestimmten Region enthält, kann bei einem zu langen Training diese Gesichter überpropotional stark gelernt werden. Wenn das KI-Modell zu lange trainiert wurde, wird dies auch als Overfitting(Überanpassung) betitelt. Das heißt, dass bestimmte Datenpunkte zu ausführlich betrachtet werden. Wenn jetzt ein overfitted Model angewandt wird auf einen neuen Datensatz, so werden eher die Gesichter bevorzugt womit das Modell trainiert wurde und so können definierte Bevölerkungsgruppen systematisch benachteiligt werden \cite{Bias}. 

Auch während der Entwicklung kann ein Bias enstehen wenn KI-Systeme genutzt werden, um Entscheidungen in verschiedenen Bereichen zu treffen, jedoch diese Entscheidung nicht genauer hinterfragt wird. Als hypothetischen Fall kann über ein KI-System bei der Polizei angenommen werden. Wenn die Polizei an einer Einsatzstelle geschickt werden, wo ein KI-System beurteitl hat, dass dieses Viertel gefährlich ist. Die schlussfolgerung wenn die Polizei an eine Einsatzstelle geschickt wird mit der Nachricht, dass die genannte Einsatzstelle gefährlichz ist, dann wrid diese auch entsprechend handeln \cite{dastin_insight_2018, Bias}.

Ein weiterer Apspekt bezüglich dem AI-Bias ist die Fairness. Wie oben beschrieben kann ein AI-Bias an jeder Stelle des Prozesses entstehen. Es existieren unterschiedliche Ansätze, um Fairness zu beurteilen. Wie zum Beispiel das AI Fariness 360 Tool von IBM oder Farilearn von Microsoft \cite{Fairness360,fairlearn}. Jedoch unterliegt der Fariness ein grundlegendes Problem. Fairness kann immer nur in Bezug auf bestimmte Gruppen oder Individuen gemessen werden und die Welt ist aufgeteilt in unterschiedlichste Gruppen und Individuen, jedoch basiert jedes aktuelle LLM auf dem WEIRD Datensatz. Der WEIRD Datensatz beschreibt die Populationsgruppe der Menschen welche an der Küste des nordöstlichen Atlantiks leben. Das heißt das das Datensazt die Menschen aus den folgenden Ländern umfasst: United States, United Kingdom, Western Europe, Australia, New Zealand, Canada. Gleichzeitig existieren Unterschiedliche Definitionen von Fariness aus verschiedenen Kulturkreisen. Das Paper "Which Humans?" zeigt, dass es unterschiedliche darstellungen von Kultur gibt in denen Fairness auch anders gewichtet wird. Es ist überhaupt schwierig zu sagen, dass es Fairness gibt als auch eine LLM zu haben, welche für alle Menschen Fairness bietet und man aktuell nur WEIRD Datensätze gibt \cite{Fairness}. 

\section{Einsatz von manipulationtechniken von künstlichen Intelligenzen}
Im folgenden würde ich gerne die Studie "How Far Did They Go? The Persuasive Tactics of Covert LLM Agents in a Discontinued Field Experiment" der Universität Zürich vorstellen. Diese Studie untersucht inwiefern LLM-Modelle Menschen in echten Foren und Dikussionen manipulieren können. Die Studie versucht die 3 grundlgenden Fragestellungen \enquote{RQ1: How does the model target or adopt social identities to establish epistemic authority and influence argumentative uptake?} \cite[S. ~2]{HowFar}. Diese Forschungsfrage untersucht, wie KI-Modelle sich eigene Identitäten aneignen, um Argumentative besser zu antworten und auch wie sehr das KI-Modell die Identatität des gegenübers berücksichtigt. \enquote{RQ2: Welche Autoritäts- und Ausrichtungsmechanismen prägen die epistemische Positionierung in von LLM generierten, auf Überzeugungsarbeit ausgerichteten Kommentaren innerhalb eines Live-Diskussionsforums?}\cite[S. ~2]{HowFar}. Diese Forschungsfrage untersucht in wie fern das KI-Modell einem zustimmt oder wiederspricht, um Recht zu haben. \enquote{RQ3: Welche kognitiven Verzerrungen werden durch diese Strategien ausgelöst oder verstärkt?}\cite[S. ~2]{HowFar}. Die dritte Forschungsfrage zielt daraufhinab zu untersuchen, in wie fern die KI-Text wiedergibt, sodass das Hirn nicht mehr Logisch sondern nach heuristiken denkt, gemäß dem Buch Think, Fast and Slow \cite{kahneman2011thinking}. 

Das grundlgende Vorgehen der Studie bestand darin, dass in den Reddit Forum r/ChangeMyView 33 unterschiedliche KI-Accounts untergeschleußt wurden. Diese haben innerhalb von 4 Moanten 1532 KI-generierte Kommentare gepostet. 

Die Auswertung erfolgte mithilfe von LLM-Auswertungsmethoden, welche durch Menschliche zwischenproben übereinstimmten.

Das Ergebnis der Studie für die erste Forschungsfrage, zeigt, dass die KI-Modell in 67,2 \% der Fälle auf Identitätsmerkamle zurückgriffen, 46,1 \% auf die auf die persönlichen Merkmale des gegenübers zurückgriffen als auch die selber zu 46,1 \% selber eigene identitäten annahmen.

Das Ergebnis der zweiten Forschungsfrage, zeigt, dass KI meistens eine negative Einstellungen in einer Konversation einnimmt. In 93,3 \% der Fällen hat versucht die KI zu wiedersprechen, aber hat auch in 72,6\% zugeständnisse abgeliefert.

Das Ergebnis der dritten Forschungsfrage, zeigt, dass die KI Verfügbarkeits-, Repräsentationheuristiken und Attributsubstitution nutzt. Bei der KI Verfügbarkeits heuristik werden emotionalgeladene Beispiele genutzt, um eine Tatsache häufiger darzustellen, als diese wirklich auftaucht. Dies ist in 43,6\% der Fällen augetreten. Bei der Repräsentationheuristik wurde versucht mit Stereotypen zu argumentieren. Dies wurde in 40,5\% der Fälle genutzt. Bei der Attributsubstitution wird versucht ein komplexes Problem auf einfache Probleme abzubilden, worin die Antwort schon bekannt ist. Dies wurde in 40,4\% der Fällen genutzt.


Diese Studie kann der datenethischen Kategorie des AI-Bias als auch dem Alignment Problem zugeordnet werdenn. Wie oben schon aufgeführt unterteilt sich der AI-Bias in die Datenerhebung, Model trainierung und Model Einsatz ein. In dieser Studie wurde ein KI-Model aktiv eingesetzt, um Interaktionen in Foren zu manipulieren, weswegen hier von einem Interaktionsbias gesprochen werden kann. Gleichzeitig wird hier das Alignment-Problem darstellt indem das Model versucht Menschen zu manipulieren und ethische bedenken ignoriert. Dabei wurde

Als Fazit ist das Ergebnis erschreckend als auch kontrovers. Zuerst wird mit dieser Studie nochmals betont, wie sehr sich LLM-Modelle in unsere gesellschft einarbeiten können und explizit genutzt werden können, um Menschen zu manipulieren in Foren. Zum anderen auch sehr grotesk, da man ab sofort davon ausgehen kann, dass die meisten Konversationen als auch Diskustionen online manipuliert sind. 

Die Studie ist auch sehr kontrovers, da eine manipulation von KI durch explizite Gesetzte des EU AI Acts existieren, welchen eine solche Nutzung eigentlich untersagen sollten, aber die Studie es geschafft hat dies unbemerkt für mehrere Monate zu tun. 

\section{Die Infokalypse und die Gefähr für demokratische Gesellschaften}

\paragraph{Infokalpyse}

(200 Wörter)
Das Verbreiten von falschinformationen war in der Vergangenheit gängige politsche praxis. Dabei wurden synthetische Inhalte generiert, welche dann verbreitet wurden. Ein bekanntes Beispiel ist die Zensur von Bildern in der Soviet Union. Dabei wurden aufwendig Bilder von Stalin bearbeitet nachdem dieser politische Gegner töten lassen hatte \cite{noauthor_censorship_2026}. In der Vergangenheits gab es ebenfalls Falschinformationen jedoch wird heutzutage von einer Infokalpyse gesprochen. Also einer Apokalypse von Informationen.
 
In der Vergangenheit ist es aufwendiger gewesen synthetische Inhalte zu ersellen. Die Infokalypse ist erst in den vergangenen Jahr verstärkt aufgetreten, da man mithilfe von KI einfacher solche Inhalte erstellen kann. Darunter zählen KI generierte Texte als auch Deepfakes. Dabei braucht es letztlich nur 2 Bilder und damit kann schon innerhalb kürzester Zeit ein Deepfake erstellt werden. Der gesamte Prozess ist einfacher geworden.

Die Verbreitung solcher Inhalte ist mit dem Internet viel einfacher geworden. Gezeigt Inhalten können innerhalb von wenigen Sekunden mehrere tausende menschen erreichen. 


Der letzte Faktor ist, dass die Inhalte jetzt von originalen Werken nicht mehr zu unterscheiden sind. Im Jahr 2023 waren KI-Generiete Inhalte noch so schlecht, dass diese direkt erkannt wurden. Heutzutage ist es für viele kaum noch möglich diese zu erkennen. 

(200 Wörter folgen)

Die allgemeinen Folgen der Infokalpyse bestehen darin zu einem die Integrität von jeglichen Informationen aufzulösen und zum anderen Menschen in Manipulative richtigungen zu lenken.

Das mit einhergehende Problem ist, dass Menschen sich das gerne aussuchen wollen, was sie glauben wollen, da jede Seite argumentative schon vorhanden ist. Wenn ich als Person Stellung zu meiner Position beziehen will, so kann ich einfach nach meiner Position suchen und mir werden Argumente geliefert, welche ich unterstütztend einsetzen kann. Jedoch ist es so, dass die gennanten Argumente auf Falsch Informationen beruhen. 


\paragraph{Folgen für eine demokratische Gesellschaft}
Diese Tatsache ist für demokratische Gesellschaften besonders gefährlich. Damit man versthet weswegen die Infokalpyse so gefährtlich für demokratische Gesellschaften ist, sollte erstmal dargestellt werden, was eine demokratische Gesellschaft ausmacht. Eine demokratische Gesellschaft basiert auf politische partizipationBürgern, welche rational handeln, diskutieren und sich einigen können über konkrete Themen. Mit der Infokalypse werden genau diese Eigenschaft geschwächt. 

(1. Argument)
Die Gesellschaft wird durch die Informationsflut manipuliert. Eins der Auffälligsten Beispiel, ist die US-Amerikanische Wahl 2016. Dabei hat Donald Trump die Daten von Cambridge Analytica genutzt, um gezielt Falschinformationen an die Bevölkerung zu verbreiten. Das Resultat ist, dass 87 Millionen Wählerprofile manipuliert wurden. Das ist mehr als die Hälfte aller Wahlberechtigten der Vereinigten Staaten. Jedoch existieren ebenfalls Beispiele welche konkret in Deutschland stattfinden. Dabei können schwurbler als auch die AFD als Beispiel genannt werden. Dabei wurden unterschiedliche Beispiele aufgeführt, sodass die Bürger an Wahlen teilgenommen haben, aber letztlich nicht diskustiert haben. 

(2. Argument)
Die Bürger handeln indirekt irrational handeln. Das meiste handeln basiert auf zugrundeliegenden prämissen. Wenn diese aber falsch sind, dass führt es dazu, dass die Entscheidung letztlich irrational ist. Wenn die Bevölkerung eben auf falschinformation agieren, wird der größte Vorteil der Gesellschaft, dass es eben gemeinsame beschlüsse gibt, welche durch diskurs und rationales denken zustande kommen, in frage gestellt.

Im gegensatz zu autoritäten Systemen besitzen 


\section{KI-Werzeuge im Bildungsprozess}
Um den Lösungsraum von KI-Werkzeugen im Bildungsprozess einschätzen zu können, wird zuerst beleuchtet, welche Vor- als auch Nachteile die Nutzung von KI-Werkzeugen abbilden im Bezug auf die Gesellschaftlicher und individueller Ebene. 

\paragraph{Einsatz von KI-Werkzeugen in Bildungsprozessen}

Das am häufigsten genannte Argument ist, dass das Denken verlernt wird. So existiert die Berühmte Studie von <>, welche belegt, dass schüler welche Aufsätze Mit KI-schreiben und daraufhin ohne KI schreiben, dass jenes Ergebnis deutlich schlechter ausgefallen ist, wenn KI eingesetzt wird. Das Argument ist dementsprechend erschreckend, weil beim Lernen nicht wirklich gelernt wird. Das weitere Problematische ist auch, dass die Leuten sich nicht mal daran erinnern können was sie geschrieben haben. X Prozent konnten nicht einmal einen einzelnen Satz wiedergeben. Es wird ein Gefühl vermittelt, dass etwas gelernt wird und dies nicht überprüft wird. 

Dabei wird nicht nur das Denken vernachlässigt, sondern das kritische auseinander setzen, also Hinterfragt, weswegen die Ergebnisse zustande kommen ebenso schlechter sind. In der akademischen Einrichtungen sollen nicht nur Informationen vermittelt werden, sondern auch Kompetenzen. Eine kompetenz ist das kritische Hinterfragen und Reflektieren von Aufgaben. So können nicht mal die Ergebnisse ihrer Qualität eingeschätzt werden, obwohl auch die Studie zeigt, dass die Verwendung von KI zu einem Prozessbeschleunigung führt, wenn man den Inhalt auch bewerten kann.

Ein weiteres Argument, was für die Nutzung von KI spricht ist das ökonomische Modell von Galor \& Tsiddon. Mit fortschreitender Technologie wird werden alte Strukturen fallen und die Denk-Muster können sich durchsetzen, welche sich anpassen können. So existiert mit neu entwickelter Technologie eine Chance die unterschied . Es ist unabdingbar, dass KI ebenfalls in der Industrie eingesetzt wird. 

Gleichzeitig sollte gesagt werden, dass es keine Frage ist, ob KI genutzt wird oder nicht, sondern eine Frage des Wie. Erste Studienzeigen, dass die Nutzung von KI im Sinne eines KI-Tutors die gleichen effekte erzielt. Also wenn KI als Lehrer eingesetzt wird, werden ebenso gute Ergebnisse erzielt. Dabei hat der Einsatz der KI einen entscheidenen Vorteil. Dieser besteht in der Verfügbarkeit und der schnelle. Manche Familien sind nicht in der Lage sich Nachhilfe für ihre Kinder zu leisten, nun existiert die Möglichkeit Hausuafgaben von einem KI-Tutor zeigen zu lassen. 

Aus der Gesellschaftlicher Sicht zeigt sich würde man bei bloßer benutzung von KI eine Gesellschaft erzeugen, welche das Denken verlennen. Wenn ich Aufgaben erfüllen, um sie zu erfüllen und nicht aus persönlichen Interesse, dann wird man auch nichts lernen. Aus individueller Sicht sind die neu entstandenn KI-Werkezuge einer der besten technologischen Fortschritte welche erzielt werden könnten. 



\paragraph{Persönlicher Einsatz von KI-Werkzeugen}

Als erste Prämisse würde ich gerne festlegen, dass künstliche Intelligenz nur ein Werkzeug ist. 
Ein Werkzeug, welches universal ist und vielseitige Probleme lösen kann. Der Prozess des Lernens beinhaltet das Lösen von unterschiedlichen Problemen. 
Auf abstrakter Ebene eben das Hinterfragen von unterschiedlichen Perspektiven, das Verstehen von Zusammenhängen, das Entwickeln von eigenständigen lösungsorientierten Denkweisen.
Auf technischer Ebene trainiert man solche Kompetenzen durch das Lösen von Aufgaben. Aus persönlichen Erfahrungen kann ich berichten, dass Heraussuchen von Informationen 
und das Korrigieren von Aufgaben entweder. 

Mit diesem Universalem Werkzeug kann ich nicht nur Probleme Lösen sondern auch Probleme erzeugen, welche ich lösen kann. 
Damit kann ich mir effektiv Zeit sparen und Probleme lösen, welche ich gerne lösen möchte.

Gleichzeitig kann ich mit KI Aufgaben delegieren, welche von keiner größeren Bedeutung sind. KI funktioniert gut im Web-Bereich. Das liegt daran, dass die Web-Programmierung 
stark repetitive Aufgaben sind. HTML und CSS Klassen zu schreiben zeugen von keiner großeren Mentalen Aufgabe, sondern 


\newpage
Vorlage Eidesstattliche Erklärung
Hiermit bestätige ich, dass ich die vorliegende Prüfung zu Wissenschaftstheorie und Ethik vom 28.06.2026 selbständig
bearbeitet und keine anderen als die erlaubten Hilfsmittel benutzt habe.
Ich erkläre hiermit explizit, dass auf den Einsatz von KI-gestützten Schreibwerkzeugen
vollständig verzichtet wurde.
Erlind Sejdiu, 28.06.2026, Karlsruhe 

\newpage
\printbibliography[heading= bibintoc, title={Literaturverzeichnis}]
\newpage
\listoffigures
\addcontentsline{toc}{section}{Abbildungsverzeichnis}
\newpage
\listoftables
\addcontentsline{toc}{section}{Tabellenverzeichnis}

\end{document}